\documentclass[a3paper,landscape,12pt]{article}
\usepackage[utf8]{inputenc}
\usepackage[spanish]{babel}
\usepackage[margin=2cm]{geometry}
\usepackage{tcolorbox}
\tcbuselibrary{skins, raster}
\usepackage{fontawesome5}
\usepackage{xcolor}
\usepackage{palatino} % Fuente elegante

% Paleta de colores profesionales
\definecolor{primary}{RGB}{30, 58, 138}    % Azul oscuro
\definecolor{secondary}{RGB}{59, 130, 246} % Azul vibrante
\definecolor{accent}{RGB}{16, 185, 129}    % Esmeralda
\definecolor{lightbg}{RGB}{248, 250, 252}  % Fondo muy claro
\definecolor{layer1}{RGB}{39, 174, 96}     % Verde (Directivas)
\definecolor{layer2}{RGB}{142, 68, 173}    % Púrpura (Orquestación)
\definecolor{layer3}{RGB}{211, 84, 0}      % Naranja oscuro (Ejecución)

\pagecolor{lightbg}
\pagestyle{empty}

\tcbset{
    base/.style={
        enhanced,
        colframe=primary,
        colback=white,
        coltitle=white,
        fonttitle=\bfseries\large,
        title={#1},
        attach boxed title to top left={xshift=5mm, yshift=-3mm},
        boxed title style={
            colback=primary,
            rounded corners,
            arc=2mm,
            boxrule=0pt
        },
        boxrule=1.2pt,
        arc=3mm,
        drop fuzzy shadow
    },
    innerbox/.style 2 args={
        enhanced,
        colframe=#1,
        colback=#1!5,
        coltitle=white,
        fonttitle=\bfseries,
        title={#2},
        boxrule=1pt,
        arc=2mm,
        drop fuzzy shadow
    }
}

\begin{document}

\begin{center}
    \vspace*{0.2cm}
    \huge\textbf{\textcolor{primary}{\faRobot\ Arquitectura del Agente IA y Marco Operativo}} \\[0.3cm]
    \Large\textcolor{secondary}{Orquestación Autónoma de 3 Capas, Memoria Evolutiva y Self-Healing}
    \vspace{0.5cm}
\end{center}

\begin{tcbraster}[raster columns=3, raster equal height=rows, raster column skip=1cm]
% Columna 1: Identidad y Rol
\begin{tcolorbox}[base=\faIdBadge\ Identidad y Entorno]
\textbf{Rol Principal:} Orquestador Autónomo
\begin{itemize}
    \item Actúa como \textbf{middleware} entre la intención humana y la ejecución determinista.
    \item Modifica archivos y ejecuta comandos en el sistema de manera autónoma.
    \item \textbf{Límite de Recursos:} 4GB RAM estrictos.
    \item \textbf{Stack de Entorno:} SO Linux, Conda (Python), Docker (C/C++ y \texttt{build-essential}).
\end{itemize}
\vspace{0.3cm}
\begin{tcolorbox}[innerbox={red!80!black}{\faExclamationTriangle\ REGLA DE ORO}]
\textbf{Nunca ejecutar ni depurar a ciegas.} Conocer el entorno real antes de operaciones complejas.
\end{tcolorbox}
\vspace{0.3cm}
\textbf{\faBrain\ Protocolo de Memoria (ChromaDB)}
\begin{itemize}
    \item \textbf{Consulta:} Buscar errores y flujos pasados.
    \item \textbf{Registro:} Guardar nuevos hallazgos y correcciones.
    \item \textbf{Autocorrección:} Apoyarse en memoria antes de soluciones desde cero.
\end{itemize}
\end{tcolorbox}% <-- IMPORTANTE: Sin saltos de línea para que sigan en la misma fila del raster
% Columna 2: Las 3 Capas
\begin{tcolorbox}[base=\faLayerGroup\ La Arquitectura de 3 Capas]
\begin{tcolorbox}[innerbox={layer1}{\faFileCode\ Layer 1: Directives (Qué hacer)}]
Manuales de operación estandarizados (SOPs).
\begin{itemize}
    \item \textbf{Formato:} Archivos \texttt{.yaml} en \texttt{directives/}.
    \item \textbf{Contiene:} Goal, Inputs, Steps, Outputs.
\end{itemize}
\end{tcolorbox}
\vspace{0.2cm}
\begin{tcolorbox}[innerbox={layer2}{\faNetworkWired\ Layer 2: Orchestration (Tú)}]
El núcleo de inteligencia (IA probabilística).
\begin{itemize}
    \item Enruta tareas y sanitiza entradas.
    \item Gestiona fallos (\textit{Self-Healing}).
    \item Delega toda lógica en scripts deterministas.
\end{itemize}
\end{tcolorbox}
\vspace{0.2cm}
\begin{tcolorbox}[innerbox={layer3}{\faCogs\ Layer 3: Execution (Cómo hacerlo)}]
Entorno reproducible y determinista.
\begin{itemize}
    \item \textbf{Herramientas:} Scripts Python o Docker.
    \item Validan sus salidas y fallan ruidosamente.
    \item \textbf{Salida:} Preferiblemente JSON o Stdout.
\end{itemize}
\end{tcolorbox}
\end{tcolorbox}% <-- IMPORTANTE: Sin saltos de línea
% Columna 3: Flujo y Self-Annealing
\begin{tcolorbox}[base=\faProjectDiagram\ Algoritmo de Ejecución]
\textbf{Flujo Paso a Paso}
\begin{enumerate}
    \item \textbf{Diagnóstico:} Extraer estado del sistema.
    \item \textbf{Búsqueda:} Leer directiva y memoria.
    \item \textbf{Pre-Análisis:} Predecir resultados.
    \item \textbf{Planificación:} Mapear herramientas.
    \item \textbf{Ejecución:} Operar terminal y archivos.
    \item \textbf{Trazabilidad:} Guardar en \texttt{run\_state.json}.
    \item \textbf{Validación:} Comparar output real y esperado.
    \item \textbf{Notificación:} Alertar vía audios.
    \item \textbf{Limpieza:} Eliminar basura temporal.
\end{enumerate}
\vspace{0.3cm}
\begin{tcolorbox}[innerbox={accent!90!black}{\faStethoscope\ Self-Annealing (Autocuración)}]
\begin{itemize}
    \item \textbf{Presupuesto:} Máximo 3 reintentos.
    \item \textbf{Análisis de Causa:} Lógica, Entorno o Recursos.
    \item \textbf{Principio:} Fiabilidad $>$ Velocidad.
\end{itemize}
\end{tcolorbox}
\end{tcolorbox}
\end{tcbraster}

\vfill
\begin{center}
    \small \textcolor{gray}{\textit{Generado automáticamente basado en AGENT\_FRAMEWORK.md y AGENT\_INSTRUCTIONS.md}}
\end{center}

\end{document}
